\scanpage{87}\mspage{31}
% L1-4
Donc il en faut rabattre ! On peut considérer $i_X$ au mieux comme une flèche de $\underline{\mathrm{Fib}}_0\,S$, et en prenant les composés on trouve \deleted{\unsure{donc}} une flèche dans $\underline{\mathrm{Fib}}_0\,S$ :
% L5-6
\diagram{\begin{tikzcd}[ampersand replacement=\&]
 X\arrow[r,"i_X"]\arrow[rr,bend right=22,no head,"{\deleted{\unsure{B}}\beta_X\ (= \alpha_X)}"']\&\Phi(X)\arrow[r,"{\Phi(j_X)}"]\&\Phi C(X),
\end{tikzcd}}
% L7-9
i.e. (regardant \inserted{maintenant} $\Phi_0$ comme foncteur à valeurs dans $\underline{\mathrm{Fib}}_0\,S$) un hom \inserted{de} foncteurs \deleted{\unsure{en $X$}}
\[
 \deleted{\unsure{B}}\beta:\mathrm{id}_{\underline{\mathrm{Fib}}_0\,S}\longrightarrow\Phi_0 p_0,
\]
% L10-13
où $p_0:\underline{\mathrm{Fib}}_0\,S\hookrightarrow\mathrm{Cat}_{/S}$ est l'inclusion canonique. \unsure{Cet|Ce} $\beta$ est dans $W$ argument par argument, et de $=$ pour $\alpha$ page 29, interprété cette fois comme
\[
 \mathrm{id}_{\mathrm{Cat}_{/S}}\longrightarrow p_0\Phi_0.
\]
% L14
Ce qu'on a prouvé finalement est ceci :

% L15-19
\noindent\begin{tabular}{|p{0.94\linewidth}@{}}
\uline{Théorème-emballage} : Soit $\unsure{\mathcal F|F}$ une sous-catégorie \inserted{pleine} de $\mathrm{Cat}_{/S}$ qui contient une des trois sous-catégories pleines $\mathrm{Fib}_0\deleted{\unsure{(}}S\deleted{\unsure{)}}$, $\underline{\mathrm{Parf}}\,S$, $\underline{\mathrm{Cofib}}_0\,S$, et \unsure{munissant|munissons} $\mathcal F$ du localiseur $W$ ``grossier'' (induit par celui de $\mathrm{Cat}_{/S}$). Alors l'inclusion
% L20-21
\[
 \mathcal F\longrightarrow\mathrm{Cat}_{/S}
\]
induit une équivalence sur les \unsure{localiseurs|localisations}.
\end{tabular}
% vertical rule beside the theorem
% archive sheet number: 86
